\labsection{4}{Unsupervised learning I: clustering and PCA}{sec:lab04-clustering-pca}

\begin{labproject}{Find structure without an answer key}
\textbf{Goal.} Choose $K$ from evidence, compare DBSCAN, and evaluate PCA through reconstruction and downstream accuracy.

\medskip\noindent\textbf{Setup.}\par
\begin{lstlisting}[style=mlpython]
import zipfile
from pathlib import Path

import joblib
import matplotlib.pyplot as plt
import numpy as np
import pandas as pd
from sklearn.cluster import DBSCAN, KMeans
from sklearn.decomposition import PCA
from sklearn.linear_model import LogisticRegression
from sklearn.metrics import accuracy_score, silhouette_score
from sklearn.model_selection import train_test_split
from sklearn.pipeline import Pipeline
from sklearn.preprocessing import StandardScaler

PROJECT_ROOT = Path.cwd().parent
DATA = PROJECT_ROOT / "all_datasets"
MODELS = PROJECT_ROOT / "models"
MODELS.mkdir(exist_ok=True)
\end{lstlisting}

\medskip\noindent\textbf{Part A --- How many clusters? Elbow and silhouette.}\par

Standardize age, income, and spending score for 200 mall customers; investigate useful segments without labels.

\begin{lstlisting}[style=mlpython]
customers = pd.read_csv(DATA / "mall_customers_dataset.csv")
print(customers.head())

features = ["Age", "Annual Income (k$)", "Spending Score (1-100)"]
X = StandardScaler().fit_transform(customers[features])
\end{lstlisting}

Try $K=2$--8. Inertia measures within-cluster squared distances and falls with $K$; silhouette measures cohesion versus separation.

\begin{lstlisting}[style=mlpython]
ks, inertias, silhouettes = [], [], []
for k in range(2, 9):
    kmeans = KMeans(n_clusters=k, n_init=20, random_state=42)
    labels = kmeans.fit_predict(X)
    ks.append(k)
    inertias.append(kmeans.inertia_)
    silhouettes.append(silhouette_score(X, labels))
    print(f"k={k}  inertia={kmeans.inertia_:.1f}  silhouette={silhouettes[-1]:.3f}")

fig, axes = plt.subplots(1, 2, figsize=(9.5, 3.4))
axes[0].plot(ks, inertias, marker="o")
axes[0].set_xlabel("Number of clusters k")
axes[0].set_ylabel("Inertia")
axes[0].set_title("Elbow: look for the bend")
axes[1].plot(ks, silhouettes, marker="o")
axes[1].set_xlabel("Number of clusters k")
axes[1].set_ylabel("Silhouette score")
axes[1].set_title("Silhouette: look for the peak")
plt.tight_layout()
plt.show()
\end{lstlisting}

\textbf{Inspect.} Inertia and silhouette suggest five or six groups; domain usefulness must guide the choice.

\medskip\noindent\textbf{Part B --- K-means versus DBSCAN.}\par

$K$-means assigns every customer. DBSCAN uses radius and neighbor count, leaving noise labeled $-1$.

\begin{lstlisting}[style=mlpython]
kmeans = KMeans(n_clusters=5, n_init=30, random_state=42)
k_labels = kmeans.fit_predict(X)
print("K-means silhouette:", round(silhouette_score(X, k_labels), 3))

dbscan = DBSCAN(eps=0.75, min_samples=6)
d_labels = dbscan.fit_predict(X)
is_clustered = d_labels != -1
n_clusters = len(set(d_labels[is_clustered]))
print("DBSCAN clusters:", n_clusters, " noise points:", int((~is_clustered).sum()))
if n_clusters >= 2:
    print("DBSCAN silhouette (clustered points only):",
          round(silhouette_score(X[is_clustered], d_labels[is_clustered]), 3))
\end{lstlisting}

Describe clusters using means of the original features.

\begin{lstlisting}[style=mlpython]
segments = customers.copy()
segments["kmeans_cluster"] = k_labels
segments["dbscan_cluster"] = d_labels
profile = segments.groupby("kmeans_cluster")[features].mean().round(1)
profile["customers"] = segments["kmeans_cluster"].value_counts().sort_index()
print(profile)
\end{lstlisting}

\begin{lstlisting}[style=mlpython]
fig, axes = plt.subplots(1, 2, figsize=(10, 3.9))
panels = [(axes[0], k_labels, "K-means, k = 5"),
          (axes[1], d_labels, "DBSCAN, eps = 0.75 (noise = -1)")]
for ax, labels, title in panels:
    scatter = ax.scatter(customers["Annual Income (k$)"], customers["Spending Score (1-100)"],
                         c=labels, cmap="tab10", alpha=0.8)
    ax.set_xlabel("Annual income (k$)")
    ax.set_ylabel("Spending score")
    ax.set_title(title)
    ax.legend(*scatter.legend_elements(), title="cluster", fontsize=8, loc="center right")
plt.tight_layout()
plt.show()
\end{lstlisting}

\textbf{Inspect.} $K$-means partitions income/spending segments; DBSCAN at this radius finds one cloud plus noise. Different grouping assumptions explain the contrast. Try a smaller \texttt{eps}.

\medskip\noindent\textbf{Part C --- PCA as compression, judged by a downstream model.}\par

Each digit has $8\times8=64$ pixels. Use the official training split; reserve test evaluation until the final comparison.

\begin{lstlisting}[style=mlpython]
archive_path = DATA / "optical_recognition_handwritten_digits.zip"
columns = [f"pixel_{i}" for i in range(64)] + ["label"]
with zipfile.ZipFile(archive_path) as archive:
    members = {Path(name).name: name for name in archive.namelist()}
    train = pd.read_csv(archive.open(members["optdigits.tra"]), header=None, names=columns)
    test = pd.read_csv(archive.open(members["optdigits.tes"]), header=None, names=columns)

X_dev = train.drop(columns="label") / 16.0     # pixel values 0-16 -> 0-1
y_dev = train["label"]
X_test = test.drop(columns="label") / 16.0
y_test = test["label"]
print("training digits:", len(X_dev), " test digits:", len(X_test))

fig, axes = plt.subplots(1, 10, figsize=(10, 1.4))
for ax, (_, row), label in zip(axes, X_dev.head(10).iterrows(), y_dev.head(10)):
    ax.imshow(row.to_numpy().reshape(8, 8), cmap="gray_r")
    ax.set_title(str(label), fontsize=9)
    ax.axis("off")
plt.suptitle("The first ten training digits (8 x 8 pixels each)", y=1.05)
plt.show()
\end{lstlisting}

Validate PCA budgets retaining 90\%, 95\%, or 99\% variance through classifier accuracy on training-file holdouts.

\begin{lstlisting}[style=mlpython]
X_fit, X_valid, y_fit, y_valid = train_test_split(
    X_dev, y_dev, test_size=0.25, random_state=42, stratify=y_dev
)

baseline = LogisticRegression(max_iter=2000).fit(X_fit, y_fit)
print("all 64 pixels: validation accuracy =",
      round(accuracy_score(y_valid, baseline.predict(X_valid)), 4))

for retained_variance in [0.90, 0.95, 0.99]:
    candidate = Pipeline([
        ("pca", PCA(n_components=retained_variance, svd_solver="full")),
        ("model", LogisticRegression(max_iter=2000)),
    ])
    candidate.fit(X_fit, y_fit)
    n_kept = candidate.named_steps["pca"].n_components_
    print(f"{retained_variance:.0%} variance: {n_kept:2d} components, "
          f"validation accuracy = {accuracy_score(y_valid, candidate.predict(X_valid)):.4f}")
\end{lstlisting}

Lock 95\%, refit both PCA and full-pixel pipelines on all training rows, then compare once on the official test file.

\begin{lstlisting}[style=mlpython]
final_model = Pipeline([
    ("pca", PCA(n_components=0.95, svd_solver="full")),
    ("model", LogisticRegression(max_iter=2000)),
])
final_model.fit(X_dev, y_dev)
full_model = LogisticRegression(max_iter=2000).fit(X_dev, y_dev)
pca = final_model.named_steps["pca"]

print("\nFINAL TEST (official test file)")
print("all 64 pixels accuracy:", round(accuracy_score(y_test, full_model.predict(X_test)), 4))
print("PCA components kept:", pca.n_components_,
      " variance retained:", round(pca.explained_variance_ratio_.sum(), 4))
print("PCA -> logistic accuracy:", round(accuracy_score(y_test, final_model.predict(X_test)), 4))
\end{lstlisting}

Reconstruct test digits from component scores and inspect cumulative explained variance.

\begin{lstlisting}[style=mlpython]
reconstructed = pca.inverse_transform(pca.transform(X_test.head(8)))

fig, axes = plt.subplots(2, 8, figsize=(9, 2.6))
for j in range(8):
    axes[0, j].imshow(X_test.iloc[j].to_numpy().reshape(8, 8), cmap="gray_r")
    axes[1, j].imshow(reconstructed[j].reshape(8, 8), cmap="gray_r")
    axes[0, j].axis("off")
    axes[1, j].axis("off")
axes[0, 0].set_title("original (64 numbers)", fontsize=8, loc="left")
axes[1, 0].set_title(f"rebuilt from {pca.n_components_} components", fontsize=8, loc="left")
plt.show()

plt.figure(figsize=(5.5, 3.4))
all_components = PCA(svd_solver="full").fit(X_dev)
plt.plot(np.arange(1, 65), np.cumsum(all_components.explained_variance_ratio_), marker=".")
plt.axhline(0.95, linestyle="--", color="grey")
plt.axvline(pca.n_components_, linestyle="--", color="grey")
plt.xlabel("Number of principal components")
plt.ylabel("Cumulative explained variance")
plt.title("Most of the picture lives in the first 30 directions")
plt.tight_layout()
plt.show()
\end{lstlisting}

\textbf{Inspect.} Reconstructions blur slightly while accuracy holds. The last 35 components carry under 5\% of variance.

Save PCA and classifier together to accept new 64-pixel images.

\begin{lstlisting}[style=mlpython]
model_path = MODELS / "digits_pca_logistic.joblib"
joblib.dump(final_model, model_path)

loaded = joblib.load(model_path)
first_digits = X_test.iloc[:5]
print("predicted:", loaded.predict(first_digits).tolist())
print("actual:   ", y_test.iloc[:5].tolist())
\end{lstlisting}

\begin{tryit}
Set DBSCAN \texttt{eps=0.5}; count clusters/noise and justify which result better serves marketing.
\end{tryit}
\end{labproject}
